% Dot batch 70; source SHA-256: 399b96deefb761499a04126ed6d513d384302a87eea424474ae59c0c4a67a2ea
% Archived source: dot_updates_20261124/batches/batch070/EGA第五卷_光滑二次型16.4至16.8_校正判别式.tex
% Original body lines: 16--64; integration map: integration/late_integration.json
\begingroup
\def\newscope{V-Q}
\def\OO{\mathcal O}
\def\Sym{\operatorname{Sym}}
\egap{16.4}
以下研究主要是为了给出16.2未覆盖情形中的二次型光滑性判据，也就是，在特征为$2$的域上考虑奇数维向量空间的情形。此时，$E^\vee$上的每个二次型$Q$都是退化的；但容易看出，它仍可能光滑，例如下文的标准型。

对每个正整数$n$，在$\mathbb Z^n$上引入\emph{标准二次型}$Q_n$，它是关于$n$个变量$X_1,\ldots,X_n$的整系数二次型。分两种情形\cmt{L18-070-01}{英译末项写成一次项$X_{2m+1}$；法文扫描右缘截断，无法辨认该项是否保留平方。二次齐次性及16.5对最后一维极化块的计算均要求平方。奇数情形$m=0$时，前面的配对和为空和。}：
\begin{align*}
\text{(a)}\ n=2m:\quad
Q_{2m}(X_1,\ldots,X_{2m})&=\sum_{i=1}^{m}X_{2i-1}X_{2i},\\
\text{(b)}\ n=2m+1:\quad
Q_{2m+1}(X_1,\ldots,X_{2m+1})&=Q_{2m}(X_1,\ldots,X_{2m})+X_{2m+1}^2 .
\end{align*}


\egas{16.5}{引理}
设$n$为奇正整数。考虑$n$个变量的普遍二次型
\[
Q(X_1,\ldots,X_n)=\sum_i a_i X_i^2+\sum_{i<j}b_{ij}X_iX_j,
\]
其中系数$a_i,b_{ij}$为不定元，故其判别式$d(Q)$是关于这些不定元的整系数多项式。这个多项式的\emph{含量}等于$2$，也就是说，它的各个系数的最大公因数为$2$。

先证明含量是$2$的倍数。这就是说，多项式的系数模$2$化简后所得的多项式恒为零。在特征$2$的域上，奇数维的二次型总是退化，因而判别式为零；将此应用于模$2$的普遍系数环的分式域，就得到所需的多项式恒等式\cmt{L18-070-02}{法文扫描仍排有否定词 \emph{pas}，但前文所要证明的恰是模$2$后恒为零，英译也是如此。两底本的“奇数次”在这里均指变量数为奇数；二次型并非奇次齐次式。仅在$\mathbb F_2$的有限个赋值上为零，并不足以推出多项式恒等为零。}。

要证明含量恰为$2$，只需计算$n$变量标准型的判别式。一方面，这个值必为含量的倍数；另一方面，在本节的极化约定下，每个配对块的判别式为$-1$，最后一维平方块的判别式为$2$。因此偶数维标准型的判别式是$(-1)^m$，奇数维的则是$2(-1)^m$。\cmt{L18-070-03}{两底本在此忽略标准型行列式的符号，写成偶数维为$1$、奇数维为$2$。用16.1的极化行列式，配对块的矩阵为$\left(\begin{smallmatrix}0&1\\1&0\end{smallmatrix}\right)$，行列式为$-1$。符号不影响含量。} 因此可以引入
\[
\widetilde d(Q)=\tfrac12 d(Q)
\]
这个关于$Q$的系数的多项式，其系数仍为整数。于是把$a_i,b_{ij}$取为任意环$A$中的元素，亦即给定一个$A$上的$n$变量二次型时，$\widetilde d$便有确定的值。

我们把$\widetilde d(Q)$称为“$n$变量普遍二次型$Q$的校正判别式多项式”；在任意环$A$上以二次型$q$的系数代入而得的值，称为“二次型$q$的校正判别式”。更一般地，从16.5直接得出下列命题。

\egas{16.6}{命题}
对于任一预概形$S$、其上有限秩局部自由模$E$以及二次型$Q\in\Gamma(S,\Sym^2 E)$，可以且只能以一种方式赋予$Q$一个$(\det E)^{\otimes2}$的截面$\widetilde d(Q)$，使之满足：
\begin{enumerate}
\renewcommand{\labelenumi}{(\alph{enumi})}
\item 与基变换相容，并且关于$E$的同构具函子性；
\item 若$E$处处为偶秩，则$\widetilde d(Q)=d(Q)$，其中$d(Q)$为$Q$的判别式；
\item 若$E$处处为奇秩，则$2\widetilde d(Q)=d(Q)$。\cmt{L18-070-04}{法文扫描在$\widetilde d(Q)$之后有一个像上标$2$的记号；英译不含这个上标。两个截面同属$(\det E)^{\otimes2}$，真正平方将落在$(\det E)^{\otimes4}$；并且16.5已经给出普遍恒等式。故正文不取平方。}
\end{enumerate}

注意：如果不同时（甚至不先）陈述关于同构的函子性，就不宜陈述与基变换相容的性质，因为$E$上的基变换本身只确定到同构为止。若把情形限于$E=\OO_S^n$，情况便不同；但这样做并不便于引用。

\egas{16.7}{定义}
$(\det E)^{\otimes2}$的截面$\widetilde d(Q)$称为二次型$Q$的\emph{校正判别式}。

\egas{16.8}{推论}
设$Q,Q'$分别是与$E,E'$对应的二次型，且$E$和$E'$各自的秩的奇偶性在$S$上恒定。在行列式线丛的自然识别下\cmt{L18-070-05}{法文写直和$\oplus$，英译误排为$\otimes$。在(b)中两底本均写系数$2$。两奇秩相加为偶秩，故左侧取普通行列式；右侧两奇秩块各有一个因子$2$，乘积是$4$。例如在$\mathbb Z$上取$Q=X^2$、$Q'=Y^2$，两块各自的校正判别式为1，而直和的判别式为4。}：
\begin{align*}
\text{(a) 若两者不都是奇秩，}\qquad
\widetilde d(Q\oplus Q')&=\widetilde d(Q)\widetilde d(Q'),\\
\text{(b) 若两者都是奇秩，}\qquad
\widetilde d(Q\oplus Q')&=4\widetilde d(Q)\widetilde d(Q').
\end{align*}

核验可直接由分块行列式得到。现在可以陈述本节的主要结论。
\par\endgroup
